%% LyX 2.4.4 created this file.  For more info, see https://www.lyx.org/.
%% Do not edit unless you really know what you are doing.
\documentclass[english,t]{beamer}
\usepackage{mathptmx}
\usepackage[T1]{fontenc}
\usepackage[latin9]{inputenc}
\usepackage{amssymb}
\usepackage{cancel}

\makeatletter
%%%%%%%%%%%%%%%%%%%%%%%%%%%%%% Textclass specific LaTeX commands.
% this default might be overridden by plain title style
\newcommand\makebeamertitle{\frame{\maketitle}}%
% (ERT) argument for the TOC
\AtBeginDocument{%
  \let\origtableofcontents=\tableofcontents
  \def\tableofcontents{\@ifnextchar[{\origtableofcontents}{\gobbletableofcontents}}
  \def\gobbletableofcontents#1{\origtableofcontents}
}

%%%%%%%%%%%%%%%%%%%%%%%%%%%%%% User specified LaTeX commands.
\usetheme{Warsaw}
% or ...

\setbeamercovered{transparent}
% or whatever (possibly just delete it)
\setbeamercovered{invisible} %makes overlays invisible


%gets rid of navigation symbols
\setbeamertemplate{navigation symbols}{}

\usepackage{multicol}

\usepackage{tikz}
\usetikzlibrary{calc}
\usetikzlibrary{plotmarks}
\usepackage{pgfplots}
\pgfplotsset{compat=newest}

\usepackage{calc}
\usepackage[nomessages]{fp}

\usepackage{advdate}    % Advancing/saving dates
\usepackage{datetime}   % Dates formatting
\usepackage{datenumber} % Counters for dates

%\usepackage{pgfpages}
%\pgfpagesuselayout{4 on 1}[a4paper, landscape, border shrink=7mm]
%\pgfpageslogicalpageoptions{1}{border code=\pgfusepath{stroke}}
%\pgfpageslogicalpageoptions{2}{border code=\pgfusepath{stroke}}
%\pgfpageslogicalpageoptions{3}{border code=\pgfusepath{stroke}}
%\pgfpageslogicalpageoptions{4}{border code=\pgfusepath{stroke}}
%%%Don't forget to add 'handout'

\makeatother

\usepackage{babel}
\begin{document}
\newcommand{\xmax}{2000}
\newcommand{\xmin}{0}
\newcommand{\xstep}{0} %must be xmin+n
\newcommand{\ymax}{500}
\newcommand{\ymin}{0}
\newcommand{\ystep}{0} %must be ymin+n

\newcounter{countEx} \setcounter{countEx}{0}
\newcounter{countProb} \setcounter{countProb}{0}
\resetcounteronoverlays{countEx} \resetcounteronoverlays{countProb}

\newcommand{\ex}[3]{
	\addtocounter{countEx}{1}
	\only<#1-#2>{\begin{exampleblock} {Example \chNum.\secNum.\arabic{countEx}.} #3	\end{exampleblock}}
}

\newcommand{\prob}[4]{
	\addtocounter{countProb}{1}
	\only<#1-#2>{\begin{exampleblock} {Problem  \chNum.\secNum.\arabic{countProb}.\,#3} #4	\end{exampleblock}}
}

\newcommand{\yline}[4]{
	(#2 - #4)/(#1 - #3)*(x - #1) + #2
}

\beamerdefaultoverlayspecification{<+->}

%%%%Chapter and Section numbers%%%%%
\newcommand{\chNum}{2}
\newcommand{\secNum}{3}
\newcommand{\secTitle}{Addition and Subtraction of Matrices}

\title[Section \chNum.\secNum: \secTitle]{Section \chNum.\secNum: \secTitle}

\subtitle{Finite Mathematics~~MTH 132~~Spring 2014}

\author{Dr.\,James Ashe}

\titlegraphic{\includegraphics[height=1.5cm]{/home/jim/JamesAshe/JCSU/images/jcsu_bull}}

\institute{Johnson C. Smith University}

\date{{}}

\makebeamertitle 
\begin{frame}{Matrix Dimensions}

\begin{columns}


\column{3.5cm}
\begin{itemize}
\item <1->[]$\begin{bmatrix}\begin{array}{rrr}
1 & 0 & -3\\
4 & -2 & 1
\end{array}\end{bmatrix}$
\end{itemize}

\column{3.5cm}
\begin{itemize}
\item <2->[]$\begin{bmatrix}\begin{array}{rrr}
1 & 0 & 0\\
0 & 1 & 0\\
0 & 0 & 1
\end{array}\end{bmatrix}$
\end{itemize}

\column{3.5cm}
\begin{itemize}
\item <3->[]$\begin{bmatrix}1\\
2\\
4\\
0
\end{bmatrix}$
\end{itemize}
\end{columns}
\end{frame}

\begin{frame}{Adding and Subtracting}

\begin{itemize}
\item <1>[]$\begin{bmatrix}\begin{array}{rrr}
1 & 4 & 0\\
0 & -1 & 3
\end{array}\end{bmatrix}+\begin{bmatrix}\begin{array}{rrr}
2 & 0 & 2\\
-2 & -1 & 2
\end{array}\end{bmatrix}$
\item <2>[]$\begin{bmatrix}\begin{array}{rrr}
1 & 4 & 0\\
0 & -1 & 3
\end{array}\end{bmatrix}-\begin{bmatrix}\begin{array}{rrr}
2 & 0 & 2\\
-2 & -1 & 2
\end{array}\end{bmatrix}$
\item <3>[]$\begin{bmatrix}\begin{array}{rrr}
1 & 4 & 0\\
0 & -1 & 3
\end{array}\end{bmatrix}+\begin{bmatrix}\begin{array}{rr}
2 & -2\\
0 & -1\\
2 & 2
\end{array}\end{bmatrix}$
\end{itemize}
\end{frame}

\begin{frame}{}

\ex{1}{}{At the beginning of an experiment three baby rats measured 5.6, 6.4, and 6 cm in length, and weighed 144, 138, and 149 g\, respectively.}
\begin{itemize}
\item [\textbf{a.}]Write a $2\times3$ matrix using this information.
\item <2->[\textbf{b.}]At the end of 2 weeks, their lengths were 10.2,
11.4, and 11.4 and weights were 196, 196, and 225 respectively. Find
a matrix representing how much the rats grew.
\item <3->[]\vspace{.5cm} \hspace{5.5cm}$\begin{bmatrix}\begin{array}{rrr}
5.6 & 6.4 & 6\\
144 & 138 & 149
\end{array}\end{bmatrix}$
\end{itemize}
\end{frame}

\end{document}
